% Generated draft: requires mathematical and source review.
\documentclass[11pt]{article}
\usepackage[margin=25mm]{geometry}
\usepackage{fontspec}
\setmainfont{Latin Modern Roman}
\usepackage{amsmath,amssymb,mathtools,unicode-math}
\setmathfont{Latin Modern Math}
\usepackage{enumitem,needspace,xurl,hyperref}
\setlength{\parindent}{0pt}
\setlength{\parskip}{5pt}
\setlength{\emergencystretch}{3em}
\newcommand{\N}{\mathbb{N}}
\newcommand{\Z}{\mathbb{Z}}
\newcommand{\Q}{\mathbb{Q}}
\newcommand{\R}{\mathbb{R}}
\newcommand{\C}{\mathbb{C}}
\newcommand{\F}{\mathbb{F}}
\newcommand{\A}{\mathbb{A}}
\newcommand{\PP}{\mathbb{P}}
\newcommand{\E}{\mathbb{E}}
\newcommand{\eps}{\varepsilon}
\newcommand{\dd}{\,\mathrm{d}}
\begin{document}
\section*{UnsolvedMath ID: 30004307}
\textbf{Source problem code:} catalogue entry
\subsection{Definitions and mathematical statement}
A quandle, written in the right-operation convention, is a set \(X\) with a binary operation \(\triangleright\) such that \(x\triangleright x=x\), each map \(R_y\colon x\mapsto x\triangleright y\) is bijective, and
\[
(x\triangleright y)\triangleright z=(x\triangleright z)\triangleright(y\triangleright z)
\qquad(x,y,z\in X).
\]
It is finite when \(|X|<\infty\). This convention does not resolve the source's unspecified meaning of ``abelian.''

For an abelian group \(B\), written additively, its torsion subgroup is
\[
\operatorname{Tor}(B)=\{b\in B:\text{there exists an integer }n\geq 1\text{ such that }nb=0\}.
\]
\textbf{Requested computation, subject to unresolved specifications.} For every finite abelian quandle \(X\), compute
\[
\operatorname{Tor}\bigl(H^2(X)\bigr).
\]
Here \(H^2(X)\) retains the source's notation for a second cohomology group. The source record does not specify the cohomology theory, coefficient group, any required coefficient action, or the definition of an abelian quandle. Consequently, the displayed request is a formulation schema rather than a uniquely specified computation. The torsion-subgroup definition above applies when the intended cohomology group is an abelian group; no coefficient choice or interpretation of ``abelian'' is imposed in this draft.
\subsection{Short English statement}
Determine the torsion in the second cohomology of every finite abelian quandle. The intended meaning of ``abelian'' and the cohomology theory and coefficients must be specified before this request has a unique mathematical meaning.
\subsection{Sources}
\begin{itemize}
\item Supplied local UnsolvedMath registry record: permanent ID \(30004307\), title ``Second-Cohomology Torsion of Finite Abelian Quandles,'' source problem code \texttt{catalogue entry}, collection ``Oberwolfach Reports Open Problems.'' No individual proposer is supplied; the proposed year is recorded as 2019.
\item Supplied report citation: \emph{Mini-Workshop: Algebraic Tools for Solving the Yang--Baxter Equation} (2020), Oberwolfach Reports, pp. 3207--3242, DOI \texttt{10.4171/owr/2019/51}. \url{https://doi.org/10.4171/owr/2019/51}. This is also the supplied literature reference labelled ``original OWR report.''
\item Supplied literature reference: Lebed--Mortier, ``abelian quandles and structure groups.'' \url{https://doi.org/10.1016/j.jpaa.2020.106474}.
\item Supplied literature reference: Lebed--Mortier, ``detailed preprint.'' \url{https://arxiv.org/abs/1908.06745}.
\item Literature assessment dated exactly \texttt{2026-08-22}: the supplied curation labels the problem partially solved and reports that Lebed--Mortier describe the torsion structurally through finite-group quotients and settle important subclasses, while a uniform explicit computation for all finite abelian quandles remains unavailable. This is supplied dated curation, not an independently verified assessment. The supplied background explicitly calls for independent verification of time-sensitive status and cited claims.
\end{itemize}
\end{document}
